% !PRINTED 135
\weberpage{135}
%<<CONT
divisor of $P$ of degree $p^{\beta} q$, where $0 \leqq \beta < \alpha$ is assumed, and let $\beta$ be taken as large as possible. Then, because of the assumed simplicity of the group $P$, the theorem of \S.\,28, 2. is applicable, according to which $P$ is isomorphic with a permutation group of $p^{\alpha-\beta}$ figures. Among the permutations of this group there are also some whose degree is divisible by $q$, and which therefore, decomposed into their cycles, must contain a cycle of $q$ figures, and from this it follows:

\weq{1}{p^{\alpha-\beta} > q.}

$\beta)$ Secondly, now let $Q$ be a divisor of $P$ of degree $p^{\alpha}$ and index $q$. Since by hypothesis $P$ is simple, the group denoted by $R$ in \S.\,30, which consists of the elements of $P$ permutable with $Q$, cannot be identical with $P$, since otherwise $Q$ would be a normal divisor of $P$, which, since $P$ is assumed to be a simple group, is after all impossible. But since $Q$ is a divisor of $R$, and the index of $Q$ with respect to $P$ is a prime, we must have $R = Q$. By II, \S.\,30 there are therefore $q$ and no more mutually distinct conjugate groups

\[ Q,\ Q_1,\ Q_2,\ \dots\ Q_{q-1}, \]

which are all of degree $p^{\alpha}$. Among these groups we take two, say $Q$, $Q_1$, which have a greatest common divisor $Q_0$ of degree $p^{\gamma}$, and we assume these two groups so chosen that $\gamma$ becomes \gesperrt{as large as possible}. Then $0 \leqq \gamma < \alpha$. If we now apply to the group $P$ the theorem of \S.\,28, 3., putting the group $Q$ for the two groups $Q$, $R$ of that theorem, then

\[ q = h_1 + h_2 + h_3 + \cdots, \]

and in it $h_1 = 1$, and the remaining summands $h_2, h_3, \dots$ are the indices of greatest common divisors of $Q$ with one each of the groups $Q_1, Q_2, \dots$. Hence $h_1, h_2, \dots$ are divisible by $p^{\alpha-\gamma}$, % [sic: Weber includes h_1 among the divisible terms although h_1 = 1.]
and it follows that

\[ q \equiv 1 \pmod{p^{\alpha-\gamma}}, \]

hence

\weq{2}{q > p^{\alpha-\gamma},}

and consequently, on account of (1):

\weq{3}{\gamma > \beta,\quad \gamma > 0.}

$\gamma)$ If now, then, $Q_0$ is the intersection of $Q$ and $Q_1$ of degree $\gamma$, % [sic: Weber prints gamma although Q_0 was defined to have degree p^gamma.]
we can by \S.\,31, III. determine a divisor $R$ of $Q$ of degree $p^{\gamma+1}$, of which $Q_0$ is a normal divisor, and $R$
%CONT>>
